\section{Minimal rational curves}\label{sec:minimal-curves}

We now assume that $X$ has Picard number one.  The first task is to
construct a smooth universal family of rational curves and determine its
relative tangent bundle along every member.  These facts will provide the
common geometric input for the two remaining pseudoindices.

All cohomology groups in the intersection calculations have rational
coefficients.  A class in $H^{2i}$ has \emph{weight} $i$.  For a smooth
morphism $a$ we write $T_a=\ker(da)$ and $-K_a=c_1(T_a)$.
Kodaira vanishing gives $H^i(X,\mathcal O_X)=0$ for $i>0$.
The exponential sequence and the Lefschetz $(1,1)$ theorem therefore give
\begin{equation}\label{eq:picard-cohomology}
  H^2(X,\mathbb Q)=\mathbb Q\,c_1(T_X).
\end{equation}
We also use that smooth complex Fano manifolds are rationally connected
and simply connected; see
\cite[Theorems~2.11 and~5.3.2]{CampanaPeternell1999}.

The \emph{pseudoindex} of $X$ is the positive integer
\[
 d=\min\{-K_X\cdot B:B\subset X\text{ is an irreducible rational curve}\}.
\]
Choose an irreducible component $V$ of the normalized space of rational
curves whose members have degree $d$.  Let
\begin{equation}\label{eq:universal-family}
 \begin{tikzcd}[column sep=large]
 U \arrow[r,"e"] \arrow[d,"\pi"'] & X\\
 V &
 \end{tikzcd}
\end{equation}
be its universal family and evaluation map.  Thus a fibre of $\pi$ is the
normalization of the corresponding curve, and its map to $X$ is
birational onto that curve.

\begin{proposition}\label{prop:minimal-family}
Let $X$ be a smooth complex Fano manifold of dimension $n$ with nef tangent
bundle, and let $V$ be a component of rational curves of minimal
anticanonical degree $d$.  Then $U$ and $V$ are smooth projective varieties,
$\dim V=n+d-3$, the map $\pi:U\to V$ is a smooth $\mathbb P^1$-fibration,
and $e:U\to X$ is smooth and surjective with connected fibres of dimension
$d-2$.  For every fibre $C$ of $\pi$,
\begin{equation}\label{eq:relative-tangent-splitting}
 T_e|_C\simeq\mathcal O_{\mathbb P^1}(-1)^{\oplus(d-2)}.
\end{equation}
\end{proposition}

\begin{proof}
The properties of $U,V,\pi,e$ are the minimal-family theorem recalled in
\cite[Proposition~1.6]{Kanemitsu2017}; that statement applies in arbitrary
dimension.  We recall why its hypotheses hold here.  For every
normalization $\varphi:\mathbb P^1\to X$, nefness implies a splitting
\[
 \varphi^*T_X=\bigoplus_{i=1}^n\mathcal O(a_i),\qquad a_i\geq0.
\]
In particular, each such map is \emph{free}, meaning that its pulled-back
tangent bundle is generated by global sections.  Minimality of $d$ rules
out a degeneration into a sum of rational curves: each nonzero component
would already have anticanonical degree at least $d$.  This is the
properness condition for the minimal family.  Freeness makes the spaces
of maps and the evaluation smooth.  Finally, the finite map in the Stein
factorization of a smooth proper evaluation is \'{e}tale; simple
connectedness of $X$ makes it an isomorphism, giving connected fibres.

It remains to prove \eqref{eq:relative-tangent-splitting}; here it is
essential to allow \emph{every} member of $V$.  Fix a fibre
$C\simeq\mathbb P^1$ and set $\varphi=e|_C$.  Let $H$ be the corresponding
component of $\operatorname{Hom}_{\mathrm{bir}}(\mathbb P^1,X)$.
Since all its maps are free, $H$ is smooth, and
\[
 T_{[\varphi]}H=H^0(\mathbb P^1,\varphi^*T_X).
\]
The construction of the normalized rational-curve space gives a
principal $\operatorname{PGL}_2$-bundle $H\to V$ and a compatible map
\[
 \lambda:\mathbb P^1\times H\longrightarrow U,
 \qquad e\circ\lambda(t,f)=f(t);
\]
see \cite[Sections~2.1--2.2]{BrionKannan2021}.
The normalization on the map space introduces no change because $H$ is
smooth.  Over $[\varphi]$, the map $\lambda$ identifies $\mathbb P^1$ with
$C$.

For $\xi\in T_{[\varphi]}H$, differentiating $\lambda$ in the constant
parameter direction $(0,\xi)$ produces a section of $T_U|_C$.
Its image under $de$ is precisely the section of $\varphi^*T_X$
represented by $\xi$.  Consequently
\begin{equation}\label{eq:sections-evaluation-surjective}
 H^0(C,T_U|_C)\longrightarrow H^0(C,\varphi^*T_X)
 \quad\text{is surjective}.
\end{equation}
The tangent sequence of $\pi$, restricted to $C$, is
\[
 0\longrightarrow\mathcal O(2)\longrightarrow T_U|_C
 \longrightarrow\mathcal O^{\oplus(n+d-3)}\longrightarrow0.
\]
It gives $h^0(T_U|_C)=n+d$ and $H^1(T_U|_C)=0$.
The nonnegative splitting of $\varphi^*T_X$ likewise gives
$h^0(\varphi^*T_X)=n+d$.  Thus the map in
\eqref{eq:sections-evaluation-surjective} is an isomorphism.
The cohomology sequence of
\[
 0\longrightarrow T_e|_C\longrightarrow T_U|_C
 \longrightarrow\varphi^*T_X\longrightarrow0
\]
now yields $H^0(C,T_e|_C)=H^1(C,T_e|_C)=0$.
Every line-bundle summand of $T_e|_C$ therefore has degree $-1$, proving
the assertion.  No immersion assumption on $\varphi$ was used.
\end{proof}

We can now isolate the cases that require further work.

\begin{proposition}\label{prop:pseudoindex-reduction}
Let $X$ be a smooth complex Fano sixfold with $\rho(X)=1$ and nef tangent
bundle.  Then either $X\simeq\mathbb P^6$, $X\simeq Q^6$, or its
pseudoindex is $4$ or $5$.
\end{proposition}

\begin{proof}
For a nonconstant map $\varphi:\mathbb P^1\to X$, the differential
$\mathcal O(2)\to\varphi^*T_X$ is nonzero.  At least one summand of the
nonnegative splitting of $\varphi^*T_X$ therefore has degree at least
$2$, so $d\geq2$.
If $d=2$, Proposition~\ref{prop:minimal-family} makes $e$ a smooth proper
map of relative dimension zero with connected fibres, hence an
isomorphism.  The map $\pi$ would then give a $\mathbb P^1$-fibration
$X\to V$ with $\dim V=5$.  The pullback of an ample divisor on $V$ is
nonzero and nef but has degree zero on the fibres.  This contradicts
$\rho(X)=1$.

For a general point $x\in X$, the \emph{variety of minimal rational
tangents} is the closure in $\mathbb P(T_x^*X)$ of the tangent lines at
$x$ to members of $V$ smooth at $x$.  The marked family $e^{-1}(x)$ has
dimension $d-2$, and the finite tangent-map theorem gives the same
dimension for this variety; see \cite[Theorem~3.4]{Kebekus2002}.
When $d=3$, it is therefore one-dimensional.  The Mok--Hwang classification
\cite[Theorem~15]{Mok2008} gives only $\mathbb P^2$, $Q^3$, and the
five-dimensional homogeneous contact manifold of type $G_2$; see also
\cite[Section~1]{Kanemitsu2017}.
None of the three resulting manifolds has dimension six.

If $d\geq6$, every rational curve $B\subset X$ satisfies
$-K_X\cdot B\geq\dim X$.  The numerical characterization of
Dedieu--H\"oring \cite[Theorem~1.3]{DedieuHoring2017} gives
$X\simeq\mathbb P^6$ or $Q^6$.
Thus the only cases left are $d=4$ and $d=5$.
\end{proof}
